\documentclass{beamer}
\usepackage[utf8]{inputenc}
\usepackage[T1]{fontenc}
\usepackage[portuguese]{babel}
\usepackage{amsmath}
\usepackage{amssymb}
\usepackage{tikz}
\usetikzlibrary{shapes.geometric, arrows.meta, calc, decorations.pathreplacing, patterns}

% Configuração do Tema
\usetheme{Madrid}
\usecolortheme{default}
\beamertemplatenavigationsymbolsempty

% Comando para o círculo vermelho
\providecommand{\circled}[1]{\mbox{\tikz[baseline=(char.base)]{\node[shape=circle,draw=red,thick,inner sep=0.5pt,text=red,font=\tiny\bfseries] (char) {#1};}}}

% Convenções visuais
\tikzset{
	eixo/.style={->, thick},
	vetorf/.style={->, thick, orange!90!black},
	vetord/.style={->, thick, blue!75!black},
	vetorv/.style={->, thick, magenta!80!black},
	vetorg/.style={->, thick, red!80!black},
	vetors/.style={->, thick, green!60!black},
	ponto/.style={circle, fill=black, inner sep=1.5pt}
}

% Informações do Título
\title{Aula 9: Teoria e Prática}
\subtitle{Energia Potencial e Conservação da Energia}
\date{11 Maio 2026}

\begin{document}
	
	% SLIDE 1
	\frame{\titlepage}
	
	
	% SLIDE 2
	\begin{frame}{Do trabalho total à energia potencial}
		\scriptsize
		
		Para uma partícula, ou um objeto tratado como partícula, já sabemos que
		$
		\textcolor{red}{\Delta K=K_f-K_i=W_{\text{tot}}.}
		$
		\textcolor{blue}{\textbf{A ideia nova:}} agora vamos separar o trabalho total em duas partes:
		\[
		\textcolor{red}{W_{\text{tot}}=W_{\text{cons}}+W_{\text{ncons}}.}
		\]
		
		\begin{columns}[T]
			\begin{column}{0.52\textwidth}
				\textcolor{blue}{\textbf{Forças conservativas}}
				
				São forças cujo trabalho pode ser representado como a diferença de uma função que depende da posição, chamada energia potencial \(U\):
				\[
				\textcolor{red}{\Delta U=U_{\text{final}}-U_{\text{inicial}}=-W_{\text{cons}}.}
				\]
				
				Exemplos: o trabalho da força gravitacional, \(W_{\text{cons}}=W_g\), ou o trabalho da força elástica, \(W_{\text{cons}}=W_k\). Então,
				\[
				\Delta U_g=-W_g
				\qquad \text{e} \qquad
				\Delta U_k=-W_k.
				\]
			\end{column}
			
			\begin{column}{0.48\textwidth}
				\textcolor{blue}{\textbf{Não conservativas}}
				
				São forças cujo trabalho $W_{ncons}$ não pode ser escrito como a diferença de uma energia potencial \(U\).
				
				Exemplo: atrito e forças externas aplicadas.
				
				\vspace{0.1cm}
			\end{column}
		\end{columns}
		
		\vspace{0.1cm}
		
		\textcolor{red}{\textbf{Mensagem:}} o trabalho não mudou. Apenas estamos a organizar melhor as contribuições. Como
		\[
		\Delta K=W_{\text{cons}}+W_{\text{ncons}},
		\qquad
		W_{\text{cons}}=-\Delta U,
		\]
		obtemos
		\[
		\textcolor{red}{\Delta K+\Delta U=W_{\text{ncons}}.}
		\]
		
	\end{frame}
	
	% SLIDE 3
	\begin{frame}{Energia potencial: ideia principal}
		\scriptsize
		
		\textcolor{blue}{\textbf{Energia potencial \(U\):}} a energia potencial é a energia associada à \textbf{configuração} de um sistema submetido à ação de uma força conservativa.
		\[
		\text{configuração} \equiv \text{arranjo dos objetos do sistema}
		\]
		
		Exemplos:
		\[
		\text{objeto--Terra: altura} \quad ; \quad
		\text{bloco--mola: compressão ou distensão}.
		\]
		
		\begin{columns}[T]
			\begin{column}{0.56\textwidth}
				Quando uma força conservativa realiza trabalho \(W_{\text{cons}}\) sobre uma partícula do sistema, a variação da energia potencial é
				
				\[
				\textcolor{red}{\Delta U=-W_{\text{cons}}.}
				\]
				
				\textcolor{blue}{\textbf{Interpretação do sinal:}}
				
				\begin{itemize}
					\item se a força conservativa realiza trabalho positivo, então \(U\) diminui:
					$
					W_{\text{cons}}>0 \Rightarrow \Delta U<0;
					$
					\item se a força conservativa realiza trabalho negativo, então \(U\) aumenta:
					$
					W_{\text{cons}}<0 \Rightarrow \Delta U>0.
					$
				\end{itemize}
			\end{column}
			
			\begin{column}{0.44\textwidth}
				\centering
				\begin{tikzpicture}[scale=0.74,>=stealth]
					\draw[eixo] (-0.2,-0.7) -- (-0.2,2.4) node[above] {$y$};
					\draw[thick] (-1.1,-0.7) -- (1.1,-0.7);
					\draw[fill=red!35, draw=red!80!black] (0.35,1.8) circle (0.18);
					\draw[vetorg] (0.35,1.55) -- (0.35,0.55) node[midway, right] {$m\vec g$};
					\draw[decorate, decoration={brace, amplitude=5pt}] (-0.55,-0.7) -- (-0.55,1.8);
					\node[left] at (-0.75,0.9) {\scriptsize altura};
					\node[align=center, font=\scriptsize] at (0.35,-1.15)
					{sistema\\ objeto--Terra};
					
					\begin{scope}[shift={(0,-2.8)}]
						\draw[pattern=north east lines] (-1.7,0.6) rectangle (-1.45,-0.6);
						\draw[thick] (-1.45,0) -- (-1.2,0)
						-- (-1.1,0.22) -- (-0.9,-0.22) -- (-0.7,0.22)
						-- (-0.5,-0.22) -- (-0.3,0.22) -- (-0.1,-0.22)
						-- (0.1,0);
						\draw[fill=gray!20] (0.1,-0.35) rectangle (0.75,0.35);
						\draw[vetors] (0.1,0.45) -- (-0.8,0.45) node[midway, above] {$\vec F_k$};
						\node[align=center, font=\scriptsize] at (0,-0.85)
						{sistema\\ bloco--mola};
					\end{scope}
				\end{tikzpicture}
			\end{column}
		\end{columns}
		
	\end{frame}
	
	% SLIDE 4
	\begin{frame}{Forças conservativas e dissipativas}
		\scriptsize
		
		\begin{columns}[T]
			\begin{column}{0.52\textwidth}
				\textcolor{blue}{\textbf{Força conservativa}}
				
				\vspace{0.2cm}
				
				Uma força é conservativa se o trabalho realizado entre dois pontos depende apenas dos pontos inicial e final, e não da trajetória.
				
				No caso da \textbf{força gravitacional}, o trabalho depende apenas das alturas inicial e final.
				
				Também podemos dizer que, em um percurso fechado, ou seja, quando a posição final é igual à posição inicial, temos:
				$
				\textcolor{red}{W_{\text{fechado}}=0.}
				$
				
				\centering
				\begin{tikzpicture}[scale=0.7,>=stealth]
					\draw[eixo] (-1.7,-0.2) -- (-1.7,2.8) node[above] {$y$};
					
					\coordinate (A) at (0,2.2);
					\coordinate (B) at (0,0.45);
					
					\draw[ponto] (A) node[right] {$A\;(y_i)$};
					\draw[ponto] (B) node[right] {$B\;(y_f)$};
					
					\draw[dashed] (-1.7,2.2) -- (A);
					\draw[dashed] (-1.7,0.45) -- (B);
					
					\draw[vetord] (A) .. controls (-0.95,1.65) and (-0.95,1.0) .. (B);
					\draw[vetord] (A) .. controls (1.0,1.75) and (1.0,0.9) .. (B);
					
					\draw[vetorg] (1.45,2.15) -- (1.45,1.15) node[midway, right] {$m\vec g$};
					
					\node[align=center, font=\scriptsize] at (0,-0.65)
					{para a gravidade, o trabalho depende\\ apenas de \(y_i\) e \(y_f\), não do caminho};
				\end{tikzpicture}
				
				Exemplos estudados: força gravitacional, força elástica.
			\end{column}
			
			\begin{column}{0.48\textwidth}
				\textcolor{blue}{\textbf{Força dissipativa}}
				
				\vspace{0.2cm}
				
				Uma força que não é conservativa é chamada de força dissipativa.
				
				\vspace{0.15cm}
				
				Exemplos indicados no material:
				\[
				\text{atrito cinético}, \qquad \text{força de arrasto}.
				\]
				
				No caso do atrito cinético, a energia cinética pode ser transformada em energia térmica:
				\[
				K \quad \Longrightarrow \quad E_t.
				\]
				
				\centering
				\begin{tikzpicture}[scale=0.8,>=stealth]
					\draw[thick] (-1.8,-0.55) -- (2.0,-0.55);
					\draw[fill=gray!20, draw=black] (-0.35,-0.55) rectangle (0.45,0.05);
					\draw[vetorv] (0.45,-0.25) -- (1.55,-0.25) node[midway, above] {$\vec v$};
					\draw[vetorf] (-0.35,-0.25) -- (-1.35,-0.25) node[midway, above] {$\vec f_{at}$};
					\node[align=center, font=\scriptsize, red!80!black] at (0,-1.25)
					{atrito: a energia térmica\\ aumenta};
				\end{tikzpicture}
			\end{column}
		\end{columns}
		
	\end{frame}
	
	% SLIDE 5
	\begin{frame}{Energia potencial gravitacional}
		\scriptsize
		
		A energia potencial gravitacional está associada a um sistema formado pela Terra e uma partícula próxima da superfície da Terra. Sistema: partícula--Terra
		
		\vspace{0.2cm}
		
		Se a partícula se desloca de uma altura \(y_i\) para uma altura \(y_f\), então
		
		\[
		\textcolor{red}{\Delta U=U_f-U_i=mg(y_f-y_i)=mg\Delta y.}
		\]
		
		Se escolhemos um ponto de referência \(y=0\) com \(U=0\), então
		$
		\textcolor{red}{U(y)=mgy.}
		$
		
		\begin{columns}[T]
			\begin{column}{0.54\textwidth}
				\textcolor{blue}{\textbf{Interpretação:}}
				
				\begin{itemize}
					\item subir significa \(\Delta y>0\):
					\[
					\Delta U>0;
					\]
					\item descer significa \(\Delta y<0\):
					\[
					\Delta U<0;
					\]
				\end{itemize}
			\end{column}
			
			\begin{column}{0.46\textwidth}
				\centering
				\begin{tikzpicture}[scale=0.75,>=stealth]
					\draw[eixo] (0,-0.3) -- (0,3.0) node[above] {$y$};
					\draw[thick] (-1.1,0) -- (1.1,0);
					\node[left] at (0,0) {$0$};
					
					\draw[fill=blue!25, draw=blue!80!black] (0.65,2.2) circle (0.18);
					\draw[vetorg] (0.65,2.0) -- (0.65,1.1) node[midway, right] {$m\vec g$};
					\draw[dashed] (0,2.2) -- (0.65,2.2);
					\node[left] at (0,2.2) {$y$};
					\node[right] at (0.9,2.2) {\scriptsize \(U=mgy\)};
					
					\draw[vetord] (-0.55,0.55) -- (-0.55,1.7) node[midway, left] {subida};
					\node[font=\scriptsize, red!80!black] at (-1.65,0.6) {$\Delta U>0$};
					
					\draw[vetord] (1.45,1.7) -- (1.45,0.55) node[midway, right] {descida};
					\node[font=\scriptsize, red!80!black] at (2.25,0.6) {$\Delta U<0$};
				\end{tikzpicture}
			\end{column}
		\end{columns}
		
	\end{frame}
	
	% SLIDE 6
	\begin{frame}{Exemplo: escolha do nível de referência}
		\scriptsize
		
		Uma preguiça de massa \(m=2{,}0\,\text{kg}\) está a \(5{,}0\,\text{m}\) acima do solo. Tomemos \(g=9{,}8\,\text{m/s}^2\). A energia potencial depende da escolha do ponto \(y=0\), mas a variação \(\Delta U\) não depende dessa escolha.
		
		\begin{columns}[T]
			\begin{column}{0.58\textwidth}
				\textcolor{blue}{\textbf{Valores de \(U=mgy\)}}
				
				\[
				\begin{array}{c|c|c}
					\text{referência }y=0 & y_{\text{preguiça}} & U=mgy \\ \hline
					\text{solo} & 5{,}0\,\text{m} & 98\,\text{J} \\
					\text{varanda a }3{,}0\,\text{m} & 2{,}0\,\text{m} & 39{,}2\,\text{J} \\
					\text{galho} & 0 & 0 \\
					1{,}0\,\text{m acima do galho} & -1{,}0\,\text{m} & -19{,}6\,\text{J}
				\end{array}
				\]
				
				\vspace{0.1cm}
				
				\textcolor{blue}{\textbf{Se a preguiça desce até o solo:}}
				
				\[
				\Delta y=-5{,}0\,\text{m}
				\]
				
				\[
				\textcolor{red}{\Delta U=mg\Delta y
					=(2{,}0)(9{,}8)(-5{,}0)
					=-98\,\text{J}.}
				\]
			\end{column}
			
			\begin{column}{0.42\textwidth}
				\centering
				\begin{tikzpicture}[scale=0.8,>=stealth]
					\draw[thick] (-1.0,-0.15) -- (1.0,-0.15);
					\node[right] at (1.5,-0.2) {\scriptsize solo};
					\draw[thick, brown!80!black] (0,0) -- (0,3.0);
					\draw[thick, brown!80!black] (-0.75,2.3) -- (0.8,2.5);
					\draw[fill=orange!60!brown, draw=black] (0.45,2.45) circle (0.16);
					\node[right] at (0.65,2.45) {\scriptsize preguiça};
					\draw[decorate, decoration={brace, amplitude=5pt}] (-0.45,0) -- (-0.45,2.45);
					\node[left] at (-0.7,1.25) {\scriptsize \(5{,}0\,\text{m}\)};
					
					\draw[dashed] (-1.1,1.47) -- (1.2,1.47);
					\node[right] at (1.2,1.47) {\scriptsize varanda};
					
					\draw[vetord] (0.95,2.35) -- (0.95,0.15) node[right] {\scriptsize descida};
					
					\node[align=center, font=\scriptsize, red!80!black] at (0,-0.75)
					{\(\Delta U\) depende apenas\\ de \(\Delta y\)};
				\end{tikzpicture}
			\end{column}
		\end{columns}
		
	\end{frame}
	
	% SLIDE 7
	\begin{frame}{Energia potencial elástica}
		\scriptsize
		
		A energia potencial elástica está associada ao estado de compressão ou distensão de um objeto elástico. Para uma mola que obedece à lei de Hooke,
		$
		\vec F_k=-k\vec x,
		$
		em que \(x=0\) corresponde ao estado relaxado:
		$
		\textcolor{red}{U=\frac12kx^2.}
		$
		Entre duas posições \(x_i\) e \(x_f\):
		\[
		\textcolor{red}{\Delta U=U_f-U_i
			=\frac12kx_f^2-\frac12kx_i^2.}
		\]
		
		\begin{columns}[T]
			\begin{column}{0.52\textwidth}
				\textcolor{blue}{\textbf{Observações:}}
				
				\begin{itemize}
					\item \(x\) mede o deslocamento em relação ao estado relaxado;
					\item \(x>0\): mola alongada;
					\item \(x<0\): mola comprimida;
					\item como aparece \(x^2\), a energia potencial elástica é positiva tanto na compressão como na distensão. O que pode ser negativo é a diferença entre o valor final e o valor inicial.
				\end{itemize}
			\end{column}
			
			\begin{column}{0.48\textwidth}
				\centering
				\begin{tikzpicture}[scale=0.62,>=stealth]
					\draw[eixo] (-2.5,-1.85) -- (2.8,-1.85) node[right] {$x$};
					\draw[dashed] (0,-2.05) -- (0,3.05);
					\node[below] at (0,-2.05) {\scriptsize \(x=0\)};
					
					\draw[domain=-2:2, smooth, variable=\x, thick, red!80!black]
					plot ({\x},{0.35*\x*\x+2.});
					\node[red!80!black] at (2.6,-0.6) {\scriptsize \(U=\frac12kx^2\)};
					
					\begin{scope}[shift={(0,1.35)}]
						\draw[pattern=north east lines] (-2.35,0.35) rectangle (-2.1,-0.35);
						\draw[thick] (-2.1,0) -- (-1.8,0)
						-- (-1.7,0.2) -- (-1.5,-0.2) -- (-1.3,0.2)
						-- (-1.1,-0.2) -- (-0.9,0.2) -- (-0.7,-0.2)
						-- (-0.5,0);
						\draw[fill=gray!20] (-0.5,-0.3) rectangle (0.0,0.3);
						\node[right] at (0.15,0) {\tiny comprimida};
					\end{scope}
					
					\begin{scope}[shift={(0,-0.10)}]
						\draw[pattern=north east lines] (-2.35,0.35) rectangle (-2.1,-0.35);
						\draw[thick] (-2.1,0) -- (-1.75,0)
						-- (-1.65,0.2) -- (-1.45,-0.2) -- (-1.25,0.2)
						-- (-1.05,-0.2) -- (-0.85,0.2) -- (-0.65,-0.2)
						-- (-0.45,0.2) -- (-0.25,-0.2) -- (0.0,0);
						\draw[fill=gray!20] (0.0,-0.3) rectangle (0.5,0.3);
						\node[right] at (0.65,0) {\tiny relaxada};
					\end{scope}
				\end{tikzpicture}
			\end{column}
		\end{columns}
		
	\end{frame}
	
	
	
	% SLIDE 8
	\begin{frame}{Exemplo: bloco e mola sem atrito}
		\scriptsize
		
		Um bloco de massa \(m=0{,}40\,\text{kg}\) move-se para a direita com velocidade \(v=0{,}50\,\text{m/s}\) e prende-se a uma mola inicialmente relaxada. A constante elástica é \(k=750\,\text{N/m}\). O bloco alonga a mola até parar momentaneamente.
		
		\[
		\text{Qual é o alongamento máximo }d?
		\]
		
		\begin{columns}[T]
			\begin{column}{0.7\textwidth}
				\textcolor{blue}{\textbf{Estado inicial}}
				
				\[
				K_i=\frac12mv^2, \qquad U_i=0.
				\]
				
				\textcolor{blue}{\textbf{Estado final}}
				
				No alongamento máximo, o bloco está momentaneamente parado:
				\[
				K_f=0, \qquad U_f=\frac12kd^2.
				\]
				
				Como não há atrito: $\Delta K + \Delta U = 0\Rightarrow
				K_i+U_i=K_f+U_f.
				$
				\[
				\frac12mv^2=\frac12kd^2
				\qquad \Rightarrow \qquad
				d=v\sqrt{\frac{m}{k}}.
				\]
				
				\[
				d=(0{,}50)\sqrt{\frac{0{,}40}{750}}
				=1{,}15\times 10^{-2}\,\text{m}
				\approx {1{,}2\,\text{cm}}.
				\]
			\end{column}
			
			\begin{column}{0.3\textwidth}
				\centering
				\begin{tikzpicture}[scale=0.75,>=stealth]
					\draw[thick] (-2.2,-0.6) -- (2.5,-0.6);
					\draw[pattern=north east lines] (-2.0,0.55) rectangle (-1.75,-0.55);
					
					\draw[thick] (-1.75,0) -- (-1.45,0)
					-- (-1.35,0.2) -- (-1.15,-0.2) -- (-0.95,0.2)
					-- (-0.75,-0.2) -- (-0.55,0.2) -- (-0.35,-0.2)
					-- (-0.15,0);
					\draw[fill=gray!20] (-0.15,-0.35) rectangle (0.55,0.35);
					\draw[vetorv] (0.55,0) -- (1.55,0) node[midway, above] {$\vec v$};
					\node[font=\scriptsize] at (0.25,-1.05) {início: \(K_i>0,\ U_i=0\)};
					
					\begin{scope}[shift={(0,-2.15)}]
						\draw[thick] (-2.2,-0.6) -- (2.5,-0.6);
						\draw[pattern=north east lines] (-2.0,0.55) rectangle (-1.75,-0.55);
						\draw[thick] (-1.75,0) -- (-1.45,0)
						-- (-1.35,0.2) -- (-1.05,-0.2) -- (-0.75,0.2)
						-- (-0.45,-0.2) -- (-0.15,0.2) -- (0.15,-0.2)
						-- (0.45,0.2) -- (0.75,-0.2) -- (1.05,0);
						\draw[fill=gray!20] (1.05,-0.35) rectangle (1.75,0.35);
						\draw[decorate, decoration={brace, mirror, amplitude=5pt}] (-0.15,-0.95) -- (1.05,-0.95);
						\node[below] at (0.45,-1.1) {$d$};
						\node[font=\scriptsize] at (0.4,-1.85) {fim: \(K_f=0,\ U_f>0\)};
					\end{scope}
				\end{tikzpicture}
			\end{column}
		\end{columns}
		
	\end{frame}
	
	% SLIDE 9
	\begin{frame}{Energia mecânica}
		\scriptsize
		
		A energia mecânica de um sistema é a soma da energia cinética com a energia potencial:
		
		\[
		\textcolor{red}{E_{\text{mec}}=K+U.}
		\]
		
		Se o sistema é isolado e apenas forças conservativas causam variações de energia, então \(W_{\text{ncons}}=0\) e a energia mecânica não varia:
		
		\[
		\textcolor{red}{\Delta E_{\text{mec}}=\Delta K+\Delta U=0.}
		\]
		
		Equivalentemente:
		$ \textcolor{red}{K_{f}+U_{f}=K_{i}+U_{i}.} $
		
		\begin{columns}[T]
			\begin{column}{0.5\textwidth}
				\textcolor{blue}{\textbf{O que pode mudar?}}
				
				\begin{itemize}
					\item \(K\) pode aumentar ou diminuir;
					\item \(U\) pode aumentar ou diminuir;
					\item mas, nessas condições:
					\[
					K+U=\text{constante}.
					\]
				\end{itemize}
				
				\vspace{0.1cm}
				
				\textcolor{blue}{\textbf{Se houver outras forças:}}
				
				Se existir trabalho de forças não incluídas em \(U\), como uma força externa aplicada e atrito, então
				\[
				\textcolor{red}{W\equiv W_{\text{ncons}}}
				\text{ e }
				\textcolor{red}{W=\Delta E_{\text{mec}}=\Delta K+\Delta U.}
				\]
			\end{column}
			
			\begin{column}{0.5\textwidth}
				\centering
				\begin{tikzpicture}[scale=0.85,>=stealth]
					\begin{scope}[shift={(0,1.15)}]
						\coordinate (O) at (0,0);
						\coordinate (L) at ({1.7*cos(235)},{1.7*sin(235)});
						\coordinate (B) at ({1.7*cos(270)},{1.7*sin(270)});
						\coordinate (R) at ({1.7*cos(305)},{1.7*sin(305)});
						
						\draw[fill=black] (O) circle (0.05);
						\draw[thick] (L) arc[start angle=235,end angle=305,radius=1.7];
						\draw[thick] (O) -- (L);
						\draw[thick] (O) -- (B);
						\draw[thick] (O) -- (R);
						\draw[fill=blue!25] (L) circle (0.13);
						\draw[fill=blue!25] (B) circle (0.13);
						\draw[fill=blue!25] (R) circle (0.13);
						\draw[vetorv] ($(B)+(-0.45,0)$) -- ($(B)+(0.45,0)$) node[midway, above] {$\vec v$};
						\node[left] at ($(L)+(-0.15,0.0)$) {\scriptsize \(U\) maior};
						\node[right] at ($(R)+(0.15,0.0)$) {\scriptsize \(U\) maior};
						\node[below] at ($(B)+(0,-0.18)$) {\scriptsize \(K\) maior};
					\end{scope}
					
					\node[align=center, font=\scriptsize, red!80!black] at (0,-2.05)
					{pêndulo sem atrito:\\ 
						\(E_{\text{mec}}\) constante;\\
						nos extremos: \(v=0\), \(U\) máxima;\\
						em baixo: \(v\) máxima, \(U\) mínima.};
				\end{tikzpicture}
			\end{column}
		\end{columns}
		
	\end{frame}
	
	% SLIDE 10
	\begin{frame}{Resolver problemas por conservação da energia mecânica}
		\scriptsize
		
		Quando a energia mecânica é conservada, podemos comparar dois instantes sem calcular os movimentos intermediários:
		$
		\textcolor{red}{K_i+U_i=K_f+U_f.}
		$
		
		\begin{columns}[T]
			\begin{column}{0.55\textwidth}
				\textcolor{blue}{\textbf{Procedimento prático}}
				
				\begin{enumerate}
					\item Escolher claramente o sistema:
					\[
					\text{objeto--Terra},\quad \text{bloco--mola},\quad \ldots
					\]
					
					\item Escolher o nível de referência para \(U\):
					\[
					U=0 \quad \text{em uma posição conveniente.}
					\]
					
					\item Identificar os dois estados:
					\[
					i=\text{inicial}, \qquad f=\text{final}.
					\]
					
					\item Escrever as energias em cada estado:
					\[
					K=\frac12mv^2, \qquad U_g=mgy, \qquad U_k=\frac12kx^2.
					\]
					
					\item Aplicar:
					$
					K_i+U_i=K_f+U_f.
					$
				\end{enumerate}
			\end{column}
			
			\begin{column}{0.45\textwidth}
				\centering
				\begin{tikzpicture}[scale=0.78,>=stealth]
					\node[draw, rounded corners, fill=blue!5, align=center] (s1) at (0,2.2)
					{\scriptsize escolher\\ sistema};
					\node[draw, rounded corners, fill=blue!5, align=center] (s2) at (0,1.1)
					{\scriptsize escolher\\ referência \(U=0\)};
					\node[draw, rounded corners, fill=blue!5, align=center] (s3) at (0,0)
					{\scriptsize escrever\\ estado \(i\)};
					\node[draw, rounded corners, fill=blue!5, align=center] (s4) at (0,-1.1)
					{\scriptsize escrever\\ estado \(f\)};
					\node[draw, rounded corners, fill=red!5, draw=red!80!black, align=center] (s5) at (0,-2.25)
					{\scriptsize \(K_i+U_i=K_f+U_f\)};
					
					\draw[->, thick] (s1) -- (s2);
					\draw[->, thick] (s2) -- (s3);
					\draw[->, thick] (s3) -- (s4);
					\draw[->, thick] (s4) -- (s5);
				\end{tikzpicture}
			\end{column}
		\end{columns}
		
	\end{frame}
	
	% SLIDE 11
	\begin{frame}{Exemplo: descida sem atrito}
		\scriptsize
		
		Um pedaço de queijo de massa \(m=2{,}0\,\text{kg}\) desce uma rampa sem atrito. A diferença de altura entre o ponto inicial e o ponto final é \(0{,}80\,\text{m}\).
		
		\[
		\text{A distância total percorrida na rampa não é necessária para calcular } \Delta U.
		\]
		
		\begin{columns}[T]
			\begin{column}{0.55\textwidth}
				\textcolor{blue}{\textbf{Variação da energia potencial}}
				
				\[
				\Delta y=y_f-y_i=-0{,}80\,\text{m}
				\]
				
				\[
				\Delta U=mg\Delta y
				=(2{,}0)(9{,}8)(-0{,}80)
				=-15{,}7\,\text{J}.
				\]
				
				A energia potencial diminui \(15{,}7\,\text{J}\). Se o queijo parte do repouso e não há atrito:
				
				\vspace{0.1cm}
				
				$
				K_i+U_i=K_f+U_f \rightarrow 
				K_f-K_i=-(U_f-U_i)=-\Delta U.
				$
				
				\vspace{0.1cm}
				
				Como \(K_i=0\):$
				K_f=15{,}7\,\text{J}.$
				
				\vspace{0.1cm}
				
				$
				\frac12mv_f^2=15{,}7
				\qquad \Rightarrow \qquad
				v_f=\sqrt{\frac{2(15{,}7)}{2{,}0}}
				\approx {4{,}0\,\text{m/s}}.
				$
			\end{column}
			
			\begin{column}{0.45\textwidth}
				\centering
				\hspace*{-0.35cm}
				\begin{tikzpicture}[scale=0.8,>=stealth]
					\draw[thick] (-1.7,1.4) -- (1.7,0);
					\draw[thick] (-1.7,0) -- (1.7,0);
					\draw[thick] (-1.7,1.4) -- (-1.7,0);
					
					\draw[fill=orange!50, draw=black] (-1.3,1.23) circle (0.13);
					\draw[fill=orange!50, draw=black] (1.1,0.25) circle (0.13);
					
					\node[left] at (-1.55,1.55) {\scriptsize \(i\)};
					\node[right] at (1.65,0) {\scriptsize \(f\)};
					
					\draw[dashed] (-2.35,1.23) -- (-1.3,1.23);
					\draw[dashed] (-2.35,0.25) -- (1.1,0.25);
					
					\draw[eixo] (-2.35,-0.1) -- (-2.35,1.75) node[above] {\scriptsize \(y\)};
					
					\node[left] at (-2.35,1.23) {\scriptsize \(y_i\)};
					\node[left] at (-2.35,0.25) {\scriptsize \(y_f\)};
					
					\draw[decorate, decoration={brace, amplitude=5pt}]
					(-2.05,0.25) -- (-2.05,1.23);
					\node[left] at (-2.25,0.74)
					{\scriptsize \(y_i-y_f=0{,}80\,\text{m}\)};
					
					\draw[vetorg] (-1.3,1.05) -- (-1.3,0.35) node[midway, right] {$m\vec g$};
					\draw[vetord] (-0.85,1.06) -- (0.35,0.56) node[midway, right] {};
					
					\node[align=center, font=\scriptsize, red!80!black] at (0,-0.75)
					{sem atrito:\\ \(U\) perdida vira \(K\)};
				\end{tikzpicture}
			\end{column}
		\end{columns}
		
	\end{frame}
	
	% SLIDE 12
	\begin{frame}{Exemplo: pêndulo sem atrito}
		\scriptsize
		
		Num pêndulo sem atrito, a energia do sistema pêndulo--Terra alterna entre energia cinética e energia potencial gravitacional.
		
		\[
		E_{\text{mec}}=K+U=\text{constante}.
		\]
		
		\begin{columns}[T]
			\begin{column}{0.55\textwidth}
				No ponto mais alto:
				$
				v_i=0 \quad \Rightarrow \quad K_i=0.
				$
				
				Se tomamos o ponto mais baixo como referência \(U_f=0\), e a diferença de altura é \(h\), então:
				
				\[
				U_i=mgh.
				\]
				
				No ponto mais baixo:
				$
				U_f=0, \qquad K_f=\frac12mv_f^2.
				$
				
				Aplicando conservação:
				\[
				K_i+U_i=K_f+U_f
				\qquad \Rightarrow \qquad
				0+mgh=\frac12mv_f^2+0.
				\]
				
				\[
				\textcolor{red}{v_f=\sqrt{2gh}.}
				\]
				
				\textcolor{blue}{\textbf{Nota:}} a massa cancela; a rapidez no ponto mais baixo depende só da diferença de altura.
			\end{column}
			
			\begin{column}{0.45\textwidth}
				\centering
				\begin{tikzpicture}[scale=0.9,>=stealth]
					\coordinate (O) at (0,1.4);
					\coordinate (L) at (-1.2,-0.4);
					\coordinate (B) at (0,-1.0);
					
					\draw[thick] (O) -- (L);
					\draw[thick] (O) -- (B);
					\draw[fill=gray!30] (O) circle (0.06);
					\draw[fill=blue!25, draw=blue!80!black] (L) circle (0.16);
					\draw[fill=blue!25, draw=blue!80!black] (B) circle (0.16);
					
					\draw[dashed] (-1.45,-1.0) -- (0.55,-1.0);
					\draw[dashed] (-1.45,-0.4) -- (-0.2,-0.4);
					\draw[decorate, decoration={brace, amplitude=5pt}] (-1.55,-1.0) -- (-1.55,-0.4);
					\node[left] at (-1.75,-0.7) {$h$};
					
					\node[left] at (-1.75,-0.45) {\scriptsize alto: \(K=0\)};
					\node[right] at (0.45,-0.95) {\scriptsize baixo: \(K\) máx.};
					
					\draw[vetorv] (-0.55,-1.0) -- (0.55,-1.0) node[midway, above] {$\vec v_f$};
					
					\node[align=center, font=\scriptsize, red!80!black] at (0,-1.75)
					{\(mgh \to \frac12mv_f^2\)};
				\end{tikzpicture}
			\end{column}
		\end{columns}
		
	\end{frame}
	
	
	% SLIDE 13
	\begin{frame}{Exemplo: mola com atrito}
		\scriptsize
		
		Um bloco de massa \(m=0{,}40\,\text{kg}\) prende-se a uma mola inicialmente relaxada com velocidade inicial \(v_i=0{,}50\,\text{m/s}\). A mola tem \(k=750\,\text{N/m}\). Com atrito, o bloco para momentaneamente quando a mola está alongada de
		$
		d=1{,}0\,\text{cm}=1{,}0\times10^{-2}\,\text{m}.
		$
		Determinar \(W_{\text{atrito}}\).
		
		\begin{columns}[T]
			\begin{column}{0.62\textwidth}
				\textcolor{blue}{\textbf{Energia no estado inicial}}
				
				\vspace{0.2cm}
				$
				K_i=\frac12mv_i^2
				=\frac12(0{,}40)(0{,}50)^2
				=0{,}050\,\text{J},
				\qquad
				U_i=0.
				$
				
				\vspace{0.2cm}
				
				\textcolor{blue}{\textbf{Energia no estado final}}
				
				\vspace{0.2cm}
				
				No alongamento máximo, o bloco está momentaneamente parado:
				$
				K_f=0.
				$
				A energia potencial elástica é
				
				\vspace{0.2cm}
				$
				U_f=\frac12kd^2
				=\frac12(750)(1{,}0\times10^{-2})^2
				=0{,}0375\,\text{J}.
				$
				
				\vspace{0.2cm}
				
				Como a força elástica já está incluída em \(U\), a única força que fica em \(W_{\text{ncons}}\) é o atrito:
				\[
				\Delta K+\Delta U=W_{\text{ncons}}=W_{\text{atrito}}.
				\]
				
				O trabalho do atrito é 
				$
				W_{\text{atrito}}=(K_f-K_i)+(U_f-U_i).
				$
				
				\[
				W_{\text{atrito}}
				=(0-0{,}050)+(0{,}0375-0)
				=-0{,}0125\,\text{J}.
				\]
				
				
			\end{column}
			
			\begin{column}{0.38\textwidth}
				\centering
				\begin{tikzpicture}[scale=0.72,>=stealth]
					\draw[thick] (-2.2,-0.6) -- (2.2,-0.6);
					\draw[pattern=north east lines] (-2.0,0.55) rectangle (-1.75,-0.55);
					
					\begin{scope}[shift={(0,0.75)}]
						\draw[thick] (-1.75,0) -- (-1.45,0)
						-- (-1.35,0.2) -- (-1.15,-0.2) -- (-0.95,0.2)
						-- (-0.75,-0.2) -- (-0.55,0.2) -- (-0.35,-0.2)
						-- (-0.15,0);
						\draw[fill=gray!20] (-0.15,-0.35) rectangle (0.55,0.35);
						\draw[vetorv] (0.55,0) -- (1.45,0) node[midway, above] {$\vec v_i$};
						\node[font=\scriptsize] at (0.0,-0.75) {início: \(x_i=0\)};
					\end{scope}
					
					\begin{scope}[shift={(0,-1.05)}]
						\draw[thick] (-1.75,0) -- (-1.45,0)
						-- (-1.35,0.2) -- (-1.05,-0.2) -- (-0.75,0.2)
						-- (-0.45,-0.2) -- (-0.15,0.2) -- (0.15,-0.2)
						-- (0.45,0.2) -- (0.75,-0.2) -- (1.05,0);
						\draw[fill=gray!20] (1.05,-0.35) rectangle (1.75,0.35);
						
						\draw[vetorf] (1.05,-0.35) -- (0.35,-0.35) node[midway, below] {$\vec f_{at}$};
						
						\draw[decorate, decoration={brace, mirror, amplitude=5pt}] (-0.15,-0.85) -- (1.05,-0.85);
						\node[below] at (0.45,-0.98) {$d$};
					\end{scope}
				\end{tikzpicture}
				
				\vspace{0.15cm}
				
				\[
				E_{\text{mec},i}=K_i+U_i
				=0{,}050\,\text{J}
				\]
				\[
				E_{\text{mec},f}=K_f+U_f
				=0{,}0375\,\text{J}
				\]
				
				\textcolor{red}{\textbf{Interpretação:}} o atrito realiza trabalho negativo; por isso, a energia mecânica diminui.
			\end{column}
		\end{columns}
		
	\end{frame}
	
	% SLIDE FINAL
	\begin{frame}{Resumo final: ideias e equações úteis}
		\scriptsize
		
		\begin{columns}[T]
			\begin{column}{0.5\textwidth}
				\textcolor{blue}{\textbf{Trabalho e energia cinética}}
				
				Para uma partícula:
				\[
				\textcolor{red}{\Delta K=K_f-K_i=W_{\text{tot}}.}
				\]
				
				\[
				W_{\text{tot}}=\sum_i W_i.
				\]
				
				\textcolor{blue}{\textbf{Separação do trabalho}}
				
				\[
				\textcolor{red}{W_{\text{tot}}=W_{\text{cons}}+W_{\text{ncons}}.}
				\]
				
				Para forças conservativas:
				\[
				\textcolor{red}{\Delta U=-W_{\text{cons}}.}
				\]
				
				Logo:
				\[
				\textcolor{red}{\Delta K+\Delta U=W_{\text{ncons}}.}
				\]
			\end{column}
			
			\begin{column}{0.5\textwidth}
				\textcolor{blue}{\textbf{Energias potenciais}}
				
				Gravitacional:
				\[
				\textcolor{red}{\Delta U_g=mg(y_f-y_i)}
				\qquad
				\textcolor{red}{U_g=mgy.}
				\]
				
				Elástica:
				\[
				\textcolor{red}{U_k=\frac12kx^2}
				\]
				
				\[
				\textcolor{red}{\Delta U_k
					=\frac12kx_f^2-\frac12kx_i^2.}
				\]
				
				\textcolor{blue}{\textbf{Energia mecânica}}
				
				\[
				\textcolor{red}{E_{\text{mec}}=K+U.}
				\]
				
				Se \(W_{\text{ncons}}=0\):
				$
				\textcolor{red}{K_i+U_i=K_f+U_f.}
				$
				
				\vspace{0.4cm}
				
				Se \(W_{\text{ncons}}\neq0\):
				$
				\textcolor{red}{W_{\text{ncons}}=\Delta E_{\text{mec}}.}
				$
			\end{column}
		\end{columns}
		
		\vspace{0.4cm}
		\centering
		\textcolor{red}{\textbf{Ideia central:}} forças conservativas podem ser descritas por \(U\); as forças não conservativas ficam como trabalho em \(W_{\text{ncons}}\).
		
	\end{frame}
	
\end{document}